%!TEX root = TheoretischePhysikC.tex

\section{Vorlesung: Spin}
Die Drehimpulsoperatoren erfüllen die folgende Vertauschungsrelationen:
\begin{equation}\label{eq:amcr}
	[\widehat{L}_1,\widehat{L}_2] = i\hbar \widehat{L}_3, \quad [\widehat{L}_2,\widehat{L}_3] = i\hbar \widehat{L}_1\quad\text{und}\quad [\widehat{L}_3,\widehat{L}_1] = i\hbar \widehat{L}_2.
\end{equation}
Es folgt, dass
\begin{equation}
	[\widehat{\mathbf{L}}^2, \widehat{L}_j]=0, \quad \text{für alle $j=1,2,3$,}
\end{equation}
wobei
\begin{equation}
	\widehat{\mathbf{L}}^2 \equiv \widehat{L}_1^2+\widehat{L}_2^2+\widehat{L}_3^2.
\end{equation}
Wir haben gesehen, dass die Bahndrehimpulsoperatoren $\widehat{\mathbf{L}} \equiv \widehat{\mathbf{r}}\times \widehat{\mathbf{p}}$ diese Kommutatoren erfüllen. Man sagt, dass $\widehat{\mathbf{r}}\times \widehat{\mathbf{p}}$ uns eine \emph{Darstellung} dieser Kommutatoren liefert. Eine wichtige Eigenschaft dieser Darstellung ist, dass den Gesamtdrehimpuls \emph{quantisiert} ist, im Sinne, dass die Eigenwerte des Gesamtdrehimpulsoperators $\widehat{\mathbf{L}}^2$ nur die Werte $\hbar^2 \ell(\ell+1)$ annehmen können, wobei $\ell = 0, \frac12, 1, \ldots$. Allerdings für die Bahndrehimpulsoperatoren $\widehat{\mathbf{r}}\times \widehat{\mathbf{p}}$ tauchen nur ganzzahlige Werte auf: $\ell = 0,1,2,\ldots$.

Zu jeder Darstellung gehört ein Hilbertraum, auf dem die Operatoren wirken; ausführlich diskutiert haben wir bisher allein $\mathcal{H}=L^2(\mathbb{R}^3)$.

Eine faszinierende Frage bleibt offen: Gibt es Quantensysteme, die halb-ganzzahlige Eigenwerte aufweisen? Um diese Frage zu bejahen benötigen wir einen Hilbertraum $\mathcal{H}$ und drei Operatoren $\widehat{S}_j:\mathcal{H}\rightarrow \mathcal{H}$, $j=1,2,3$, die die Relationen (\ref{eq:amcr}) erfüllen:
\begin{equation}\label{eq:amcr-spin}
	\eqbox{[\widehat{S}_1,\widehat{S}_2] = i\hbar \widehat{S}_3, \quad [\widehat{S}_2,\widehat{S}_3] = i\hbar \widehat{S}_1\quad\text{und}\quad [\widehat{S}_3,\widehat{S}_1] = i\hbar \widehat{S}_2.}
\end{equation}
Falls ein solcher Hilbertraum und solche Operatoren existieren, können wir zumindest feststellen, dass es keine \emph{mathematische} Hindernisse gibt die die Existenz von Quantensystemen mit halb-ganzzahligem gesamten Drehimpuls verhindern. Selbst wenn ein solches Quantensystem mathematisch existieren kann, ist es noch ein weiterer Schritt, festzustellen, ob ein solches System in der Natur vorkommen kann.

Es zeigt sich, dass Quantensysteme mit halbzahligen Eigenwerten für den gesamten Drehimpulsoperator tatsächlich existieren. Wir untersuchen nun den Fall $\ell = \frac12$. Der dazugehörige Hilbertraum ist durch $\mathcal{H} \equiv \mathbb{C}^2$ gegeben. Dass wir die Operatoren jetzt $\widehat{S}_j$ statt $\widehat{L}_j$ nennen, ist eine reine Umbenennung. Wir definieren dazu die Matrizen
\begin{equation}
	\widehat{S}_1 = \frac{\hbar}{2} \begin{pmatrix}
		0 & 1 \\ 1 & 0
	\end{pmatrix}, \quad \widehat{S}_2 = \frac{\hbar}{2} \begin{pmatrix}
		0 & -i \\ i & 0
	\end{pmatrix} \quad \text{und} \quad \widehat{S}_3 = \frac{\hbar}{2} \begin{pmatrix}
		1 & 0 \\ 0 & -1
	\end{pmatrix}.
\end{equation}
Bis auf den Faktor $\frac{\hbar}{2}$ sind das die \emph{Pauli-Matrizen} $\sigma_1,\sigma_2,\sigma_3$. Zu prüfen bleibt, ob sie (\ref{eq:amcr-spin}) erfüllen:
\begin{equation}
	[\widehat{S}_1,\widehat{S}_2] = \frac{\hbar^2}{4}\left(\sigma_1\sigma_2-\sigma_2\sigma_1\right) = \frac{\hbar^2}{4}\cdot 2i \begin{pmatrix}
		1 & 0 \\ 0 & -1
	\end{pmatrix} = i\hbar\, \widehat{S}_3.
\end{equation}
Die beiden übrigen Relationen folgen nach demselben Muster.

\begin{historical}
	Damit ist erst die \emph{mathematische} Frage beantwortet: Man kann wilde Hilberträume definieren, die alle Bedingungen erfüllen und trotzdem kein System in der Natur beschreiben. Beim Spin geht es gut -- im Stern-Gerlach-Experiment, das einen Silberatomstrahl in genau zwei Teilstrahlen aufspaltete, wurde eines gefunden: Das Elektron besitzt einen \emph{internen} Freiheitsgrad mit Gesamtdrehimpuls $\frac12$.
\end{historical}

Der Gesamtdrehimpulsoperator $\widehat{\mathbf{S}}^2$ lautet
\begin{equation}
	\widehat{\mathbf{S}}^2 = \frac{3\hbar^2}{4} \mathbb{I}.
\end{equation}
Die Rechnung passt in eine Zeile, denn $\sigma_1^2 = \sigma_2^2 = \sigma_3^2 = \mathbb{I}$. Dieser Operator besitzt zwei (entartete) Eigenwerte, nämlich, $\frac{3\hbar^2}{4}$ und $\frac{3\hbar^2}{4}$. Vergleichen wir die Wirkung von $\widehat{\mathbf{S}}^2$ mit
\begin{equation}
	\widehat{\mathbf{S}}^2|\ell,m\rangle = \hbar^2\ell(\ell+1)|\ell,m\rangle
\end{equation}
stellen wir fest, dass $\ell=\frac12$ sein muss, denn $\frac34 = \frac12\left(\frac12+1\right)$.

Dass $[\widehat{\mathbf{S}}^2,\widehat{S}_j]=0$ gilt, ist hier trivial: Ein Vielfaches des Einsoperators vertauscht mit allem. Bezüglich dieser Darstellung ist $\widehat{S}_3$ bereits diagonal. Das heisst, die Eigenbasis von $\widehat{S}_3$ darf mit $|\ell,m\rangle$ identifiziert werden:
\begin{equation}
	|\ell=\tfrac12,m=\tfrac12\rangle \equiv \begin{pmatrix}
		1 \\ 0
	\end{pmatrix}\quad\text{und}\quad 
	|\ell=\tfrac12,m=-\tfrac12\rangle \equiv\begin{pmatrix}
		0 \\ 1
	\end{pmatrix}.
\end{equation}
Man schreibt dafür auch kurz $|\uparrow\rangle$ und $|\downarrow\rangle$.

Darstellungen zu $\ell=\frac32,\frac52,\ldots$ lassen sich ebenso angeben, \emph{fundamentale} Teilchen mit $\ell=\frac32$ hat aber noch niemand entdeckt. \emph{Aufbauen} lassen sich solche Drehimpulse allerdings: Ein System mit $\ell=\frac12$ und eines mit $\ell=1$ zusammengebracht sieht von fern aus wie $\ell=\frac32$. Damit sind wir bei zusammengesetzten Systemen.

\section*{Zusammengesetzte Systeme}
Oft müssen wir bei der Modellierung eines quantenmechanischen Systems das System aus mehreren Teilsystemen aufbauen: aus einem Elektron und einem Proton wird dann ein Wasserstoff-Atom und so weiter; oder wir verfeinern unsere Beschreibung durch Hinzunahme weiterer Freiheitsgrade: Aus einem ``Elektron ohne Spin'', wird zusammen mit einer Beschreibung eines isolierten Spinfreiheitsgrads eine Theorie für Elektronen, die den Spin einschließt.

In all diesen Fällen verfahren wir nach der Regel, dass jeder Sorte quantenmechanischer
Systeme ein Hilbertraum zugeordnet wird. Sei also $\mathcal{H}_1$ der Hilbertraum der Systemsorte ``1'' und ebenso $\mathcal{H}_2$. Die Frage lautet also: Wie und in welchem Hilbertraum ``$\mathcal{H}_{\text{$1$\&$2$}}$'' sollen wir die Systemsorte ``1 \& 2'' von Systemen beschreiben, die aus einem Paar von Teilchen/Freiheitsgraden der Sorten 1 und 2 besteht? Die Antwort ist mit Hilfe durch das \emph{Tensorprodukt} $\mathcal{H}_1\otimes \mathcal{H}_2$ von $\mathcal{H}_1$ und $\mathcal{H}_2$ gegeben.

Aus zwei Hilberträumen lässt sich auf mehrere Weisen ein dritter bauen; nur eine davon ist die physikalisch richtige. Die Forderung lautet: Die Teilsysteme müssen sich \emph{unabhängig} voneinander präparieren lassen -- ob neben dem einen Wasserstoffatom noch ein zweites steht, darf keine Rolle spielen. Also muss jede Kombination von Zuständen einem Vektor in $\mathcal{H}_{\text{$1$\&$2$}}$ entsprechen, und daraus folgt
\begin{equation}
	\dim\left(\mathcal{H}_{\text{$1$\&$2$}}\right) = \dim(\mathcal{H}_1)\cdot \dim(\mathcal{H}_2):
\end{equation}
Die Dimensionen \emph{multiplizieren} sich, sie addieren sich nicht -- wie bei klassischen Bits, wo drei Bits acht Zustände besitzen.

\begin{center}
\parbox[c][5cm][c]{14cm}{\begin{tcolorbox}[colback=banana,arc=0cm,frame hidden,enhanced]{Das Tensorprodukt $\mathcal{H}_1\otimes \mathcal{H}_2$ zweier Hilberträume $\mathcal{H}_1$ und $\mathcal{H}_2$ ist \emph{der} Hilbertraum, der Vektoren $\phi_1\otimes \phi_2$ mit $\phi_1\in \mathcal{H}_1$ und $\phi_2\in \mathcal{H}_2$ enthält, und von ihnen erzeugt wird, und in dem die Rechenregel 
\begin{equation}
	\langle \phi_1\otimes \phi_2|\psi_1\otimes\psi_2\rangle \equiv \langle\phi_1| \psi_1\rangle\langle\phi_2| \psi_2\rangle
\end{equation}
für Skalarprodukte gilt.
}\end{tcolorbox}}
\end{center}

Betrachten wir also als nächstes das Tensorprodukt von Operatoren. Für $A_1:\mathcal{H}_1 \rightarrow \mathcal{H}_1$ und $A_2:\mathcal{H}_2 \rightarrow \mathcal{H}_2$ definieren wir den Operator $A_1\otimes A_2:\mathcal{H}_1 \otimes\mathcal{H}_2 \rightarrow \mathcal{H}_1 \otimes\mathcal{H}_2$ durch die Gleichung
\begin{equation}
	\eqbox{(A_1\otimes A_2)(\phi_1\otimes \phi_2) \equiv (A_1\phi_1)\otimes (A_2\phi_2).}
\end{equation}
Da $A_1\otimes A_2$ ja ein linearer Operator sein soll, ist er damit auch auf allen
Linearkombinationen von Produktvektoren definiert.

Ein Tensorprodukt von Operatoren lässt sich natürlich auch durch seine Matrixelemente angeben. Direkt aus der Definition folgt
\begin{equation}
	\begin{split}
			[A_1\otimes A_2]_{jj',kk'} &\equiv \langle \phi_j\otimes \phi_{j'}|\left(A_1\otimes A_2\right) \phi_k\otimes\phi_{k'}\rangle \\
			&= \langle \phi_j\otimes \phi_{j'}| (A_1\phi_k)\otimes (A_2\phi_{k'})\rangle \\
			&= \langle \phi_j|A_1\phi_k\rangle \langle \phi_{j'}|A_2\phi_{k'}\rangle \\
			&= [A_1]_{j,k}[A_2]_{j',k'}.
	\end{split}
\end{equation}
Die Indizes an der Produktmatrix, sind also Paare von Indizes. Das Tensorprodukt einer $3\times 3$-Matrix mit einer $7\times 7$-Matrix ergibt also eine $21\times 21$-Matrix. Zur numerischen Verarbeitung solcher Produkte ist es nötig, aus den Paaren von Indizes wieder einfache Indizes zu machen, also alle möglichen Paare irgendwie durchzunummerieren.

Rechenregeln für das Tensorprodukt von Operatoren folgen meist leichter aus der Definition. Wir bekommen neben der Linearität in jedem Faktor, die für alle Tensorprodukte (auch von Vektoren) gilt:
\begin{align}
	(A_1\otimes A_2)(B_1\otimes B_2) &= A_1B_1\otimes A_2B_2 \\
	(A_1\otimes A_2)^\dag &= A_1^\dag \otimes A_2^\dag \\
	\text{$A_1\ge 0$ und $A_2\ge 0$} &\Rightarrow A_1\otimes A_2 \ge 0 \\
	\|A_1\otimes A_2\| &= \|A_1\|\| A_2\| \\
	\tr(A_1\otimes A_2) &= \tr(A_1)\tr(A_2).
\end{align}

Schließlich sollte noch erwähnt werden, dass alles, was hier über das Tensorprodukt aus zwei Faktoren gesagt wurde, sofort auf endlich viele Faktoren übertragbar ist.

\section*{Dynamik zusammengesetzter Systeme}
Laufen beide Teilsysteme unabhängig voneinander, so ist die Zeitentwicklung des Gesamtsystems durch $U_t = U^{(1)}_t\otimes U^{(2)}_t$ gegeben. Beim Hamiltonoperator lauert dagegen eine Falle: Man würde vermuten, er sei das Produkt $\widehat{H}_1\otimes\widehat{H}_2$ -- das ist falsch. Richtig ist
\begin{equation}\label{eq:H-tensor}
	\eqbox{\widehat{H} = \widehat{H}_1\otimes \mathbb{I} + \mathbb{I}\otimes \widehat{H}_2.}
\end{equation}
Der Grund ist die Produktregel: Ableiten von $U_t$ nach der Zeit gibt
\begin{equation}
	\frac{d}{dt}\left(U^{(1)}_t\otimes U^{(2)}_t\right) = \left(\frac{dU^{(1)}_t}{dt}\right)\otimes U^{(2)}_t + U^{(1)}_t\otimes\left(\frac{dU^{(2)}_t}{dt}\right),
\end{equation}
und daraus liest man mit $i\hbar\frac{dU^{(j)}_t}{dt} = \widehat{H}_j U^{(j)}_t$ und $i\hbar\frac{dU_t}{dt} = \widehat{H}U_t$ genau (\ref{eq:H-tensor}) ab. Energien addieren sich, sie multiplizieren sich nicht. Mit gelösten Teilproblemen $\widehat{H}_1|E_j\rangle = E_j|E_j\rangle$ und $\widehat{H}_2|E_k\rangle = E_k|E_k\rangle$ ist auch $\widehat{H}$ diagonalisiert:
\begin{equation}
	\begin{split}
		\widehat{H}\left(|E_j\rangle\otimes |E_k\rangle\right) &= \left(\widehat{H}_1|E_j\rangle\right)\otimes |E_k\rangle + |E_j\rangle \otimes \left(\widehat{H}_2|E_k\rangle\right) \\
		&= (E_j+E_k)\, |E_j\rangle\otimes |E_k\rangle.
	\end{split}
\end{equation}
\begin{übung}
	Zeigen Sie, dass diese Produktvektoren eine \emph{vollständige} Eigenbasis von $\widehat{H}$ bilden.
\end{übung}

\section*{Verschränkung}
Eine letzte Bemerkung, und sie ist die folgenreichste. Für $\mathcal{H} = \mathbb{C}^2\otimes\mathbb{C}^2$ mit Orthonormalbasen $\{\psi_1,\psi_2\}$, $\{\phi_1,\phi_2\}$ ist jeder reine Zustand eine Linearkombination der vier Produktvektoren $\psi_j\otimes\phi_k$ -- aber nicht jede ist selbst ein Produkt. Ein Beispiel ist der Einstein-Podolsky-Rosen-Zustand ($\psi_1=\phi_1=|\uparrow\rangle$, $\psi_2=\phi_2=|\downarrow\rangle$):
\begin{equation}
	|\text{EPR}\rangle = \frac{1}{\sqrt{2}}\left(\psi_1\otimes\phi_1 + \psi_2\otimes\phi_2\right) = \frac{1}{\sqrt{2}}\left(|\uparrow\uparrow\rangle + |\downarrow\downarrow\rangle\right).
\end{equation}
Er lässt sich nicht als $\psi\otimes\phi$ schreiben: Ein solches Produkt $(a_1\psi_1+a_2\psi_2)\otimes(b_1\phi_1+b_2\phi_2)$ hat die Koeffizienten $a_jb_k$, und aus $a_1b_2 = a_2b_1 = 0$ folgt $a_1b_1 = 0$ oder $a_2b_2 = 0$ -- beide können nicht $\tfrac{1}{\sqrt{2}}$ sein. Also lässt er sich nicht als unabhängige Präparation zweier Systeme auffassen -- obwohl es gerade die unabhängige Präparierbarkeit war, die uns zum Tensorprodukt geführt hat. Das Tensorprodukt liefert also mehr, als wir von ihm verlangt haben.

% fig_vl17_bloch.tex — Spin 1/2: die Blochkugel eines einzelnen Spins, und
% was aus dem Blochpfeil wird, sobald der Spin zu einem verschränkten Paar
% gehört.
%
% Follows the conventions of fig_vl11_kasten (the reference plate).
%
%   * geometry comes from the formulae of the lecture, never freehand.
%
%     Panel (a).  Every pure state of H = C^2 is, up to a global phase,
%         |psi(theta,phi)> = cos(theta/2)|up> + e^{i phi} sin(theta/2)|down>,
%     and with the S_j of the lecture its vector of expectation values is
%         (2/hbar)(<S_1>,<S_2>,<S_3>)
%             = (sin th cos ph, sin th sin ph, cos th),
%     a unit vector.  The panel therefore draws the unit sphere of that
%     vector, and every marked point sits where the formula puts it:
%     |up>,|down> at th = 0,pi — antipodal, not perpendicular, because the
%     half angle th/2 turns the 90 degrees of orthogonality in C^2 into
%     180 degrees in R^3; |+>,|-> at th = pi/2, ph = 0,pi;  |+i>,|-i> at
%     th = pi/2, ph = +-pi/2.  The projection is a genuine orthographic
%     one, camera azimuth psi = 25 deg and elevation eps with sin eps = 0.3:
%         x^ -> (-0.42262, -0.27189)
%         y^ -> ( 0.90631, -0.12679)
%         z^ -> ( 0.00000,  0.95394)
%     Consequences that are therefore forced rather than drawn by eye: the
%     equator projects to an ellipse of axis ratio 0.3; its *near* half is
%     the lower one, hence solid, the far half dashed; and the poles land
%     at 0.954 R, just inside the outline, which is what an orthographic
%     projection of a sphere actually does.  \tpcbx/\tpcby implement the
%     projection, so no coordinate in the panel was typed by hand.
%
%     Panel (b).  The one-parameter family
%         |psi_a> = cos(a)|up up> + sin(a)|down down>
%     runs from a product state through the EPR state of the lecture
%     (a = pi/4) to another product state.  Both curves are exact:
%         (2/hbar)|<S (x) 1>| = |cos^2 a - sin^2 a| = |cos 2a|,
%             the length of spin 1's arrow in panel (a) — for this family
%             <S_1 (x) 1> = <S_2 (x) 1> = 0 and <S_3 (x) 1> = (hbar/2)cos 2a,
%             all three computable with the tensor-product rules of the
%             lecture alone;
%         2|c_11 c_22 - c_12 c_21| = 2|cos a sin a| = |sin 2a|,
%             exactly the obstruction the lecture computes: a product needs
%             a_1 b_2 = a_2 b_1 = 0, i.e. a vanishing determinant of the
%             coefficient matrix.
%     Since cos^2 + sin^2 = 1 the two readings trade off exactly, and the
%     crossing points at a = pi/8, 3pi/8 are at 1/sqrt(2).  At the EPR
%     point the arrow of panel (a) has length zero: it has retreated from
%     the surface of the ball to its centre.  That is the whole content of
%     the plate — the two panels are one picture, read twice.
%
%   * colour carries the observable and nothing else: eigenstates of S_3
%     navy, of S_1 blue, of S_2 teal (the ordered ramp of `tpclevels`);
%     the one state that is not a basis state, the general |psi>, is
%     orange, and so is the entanglement curve in (b) — the blue curve in
%     (b) is the length of a blue-and-navy object, the arrow.  Axes,
%     equator and the EPR rule are neutral grey: reference structure.
%   * near/far is coded by fill, not by colour: a point of the near
%     hemisphere is a filled disc, one of the far hemisphere an open
%     circle, and the far halves of the axes are dashed.  This survives
%     greyscale printing and makes the antipodal pairs read as pairs.
%   * (b) separates its curves by dash pattern as well as by colour, for
%     the same reason.
%
% Panel (a) the state space of one spin 1/2.
% Panel (b) what becomes of that arrow when the spin is one of a pair.

% Orthographic projection of the sphere of radius #1, see the header.
% #2 = polar angle theta, #3 = azimuth phi, both in degrees.
\def\tpcbx#1#2#3{#1*(-0.42262*sin(#2)*cos(#3)+0.90631*sin(#2)*sin(#3))}
\def\tpcby#1#2#3{#1*(-0.27189*sin(#2)*cos(#3)-0.12679*sin(#2)*sin(#3)%
                     +0.95394*cos(#2))}
% \tpcbc{name}{radius}{theta}{phi} — name a point on that sphere.
\newcommand{\tpcbc}[4]{%
  \coordinate (#1) at ({\tpcbx{#2}{#3}{#4}},{\tpcby{#2}{#3}{#4}});}
\def\tpcRb{1.7}

\begin{center}
\begin{tikzpicture}[
    % near hemisphere: filled disc;  far hemisphere: open circle
    sdot/.style ={circle, inner sep=0pt, minimum size=3.4pt,
                  fill=#1, draw=#1},
    sdot/.default=cgblue,
    sdotb/.style={circle, inner sep=0pt, minimum size=3.4pt,
                  fill=white, draw=#1, line width=0.65pt},
    sdotb/.default=cgblue,
    blochaxis/.style={-{Stealth[length=5pt,width=3.5pt]}, birrengray,
                      line width=0.7pt},
    blochback/.style={birrengray, line width=0.7pt, densely dashed},
  ]

% ─────────────────────────────────────────────────────────────
%  (a)  The Bloch sphere of one spin 1/2
% ─────────────────────────────────────────────────────────────
\node[font=\sf\small\bfseries, text=cgblue] at (0,3.00)
  {(a) \tpcifenglish{one spin: the Bloch sphere}{ein Spin: die Blochkugel}};

% Every named point is placed by the parametrisation, nothing by eye.
\coordinate (O) at (0,0);
\tpcbc{Sup}  {\tpcRb}{0}  {0}      % |up>      = eigenstate of S_3, +
\tpcbc{Sdn}  {\tpcRb}{180}{0}      % |down>    = eigenstate of S_3, -
\tpcbc{Splu} {\tpcRb}{90} {0}      % |+>       = eigenstate of S_1, +
\tpcbc{Smin} {\tpcRb}{90} {180}    % |->       = eigenstate of S_1, -
\tpcbc{Spi}  {\tpcRb}{90} {90}     % |+i>      = eigenstate of S_2, +
\tpcbc{Smi}  {\tpcRb}{90} {270}    % |-i>      = eigenstate of S_2, -
\tpcbc{Spsi} {\tpcRb}{55} {-55}    % a general |psi(theta,phi)>
\tpcbc{Sfoo} {\tpcRb*sin(55)}{90}{-55}   % its foot in the equatorial plane
% Axis tips, all three at the same distance on the page (1.95 cm).
\tpcbc{Axa}  {3.880}{90}{0}
\tpcbc{Axb}  {2.131}{90}{90}
\tpcbc{Axc}  {2.044}{0} {0}
% Far ends of the three axes, just past the far surface.
\tpcbc{Bxa}  {1.05*\tpcRb}{90} {180}
\tpcbc{Bxb}  {1.05*\tpcRb}{90} {270}
\tpcbc{Bxc}  {1.05*\tpcRb}{180}{0}

% Far halves of the axes: dashed, they run behind the ball
\draw[blochback] (O) -- (Bxa);
\draw[blochback] (O) -- (Bxb);
\draw[blochback] (O) -- (Bxc);

% Equator: far (upper) half dashed, near (lower) half solid.  The axis
% ratio 0.3 = sin(eps) is forced by the projection, not chosen.
\draw[birrengray, line width=0.55pt, densely dashed]
  (-\tpcRb,0) arc (180:0:{\tpcRb} and {0.3*\tpcRb});
\draw[birrengray, line width=0.55pt]
  (-\tpcRb,0) arc (180:360:{\tpcRb} and {0.3*\tpcRb});

% Outline
\draw[birrengray, line width=0.9pt] (0,0) circle (\tpcRb);

% Near halves of the axes
\draw[blochaxis] (O) -- (Axa);
\draw[blochaxis] (O) -- (Axb);
\draw[blochaxis] (O) -- (Axc);
\node[lbl, anchor=north east, xshift=-1pt] at (Axa) {$\widehat{S}_1$};
\node[lbl, anchor=west, xshift=2pt]        at (Axb) {$\widehat{S}_2$};
\node[lbl, anchor=south, yshift=1pt]       at (Axc) {$\widehat{S}_3$};

% The general state: its arrow, its foot in the equatorial plane, and the
% two angles of the parametrisation.
\draw[tpc thin, birrengray, densely dotted] (O)    -- (Sfoo);
\draw[tpc thin, birrengray, densely dotted] (Sfoo) -- (Spsi);
\draw[-{Stealth[length=6pt,width=4pt]}, birrenorange, line width=1.1pt]
  (O) -- (Spsi);

\draw[birrengray, line width=0.55pt]
  plot[domain=0:55, samples=30, variable=\s]
    ({\tpcbx{0.62}{\s}{-55}},{\tpcby{0.62}{\s}{-55}});
\node[lbl, font=\scriptsize] at (-0.37,0.66) {$\theta$};

\draw[birrengray, line width=0.55pt]
  plot[domain=0:-55, samples=30, variable=\p]
    ({\tpcbx{0.55}{90}{\p}},{\tpcby{0.55}{90}{\p}});
\node[lbl, font=\scriptsize] at (-0.70,-0.25) {$\varphi$};

% The six eigenstates of S_1, S_2, S_3.  Colour = observable;
% filled = near hemisphere, open = far hemisphere.
\node[sdot=spacecadet]  at (Sup)  {};
\node[sdotb=spacecadet] at (Sdn)  {};
\node[sdot=cgblue]      at (Splu) {};
\node[sdotb=cgblue]     at (Smin) {};
\node[sdot=munsell]     at (Spi)  {};
\node[sdotb=munsell]    at (Smi)  {};

\node[lbl, anchor=west, xshift=3pt]       at (Sup)  {$\ket{\uparrow}$};
\node[lbl, anchor=west, xshift=3pt]       at (Sdn)  {$\ket{\downarrow}$};
\node[lbl, anchor=east, xshift=-4pt]      at (Splu) {$\ket{+}$};
\node[lbl, anchor=south west, xshift=2pt] at (Smin) {$\ket{-}$};
\node[lbl, font=\scriptsize, text=birrengray, anchor=north, yshift=-4pt]
                                          at (Spi)  {$\ket{{+}i}$};
\node[lbl, font=\scriptsize, text=birrengray, anchor=south, yshift=4pt]
                                          at (Smi)  {$\ket{{-}i}$};
\node[lbl, text=birrenorange, anchor=south east, yshift=1pt]
                                          at (Spsi) {$\ket{\psi}$};

% The formula the whole panel is drawn from, and how to read it.
\node[lbl, font=\scriptsize] at (0,-2.38)
  {$\ket{\psi}=\cos\tfrac{\theta}{2}\,\ket{\uparrow}
    +e^{i\varphi}\sin\tfrac{\theta}{2}\,\ket{\downarrow}$};
\node[lbl, font=\scriptsize, text=birrengray] at (0,-2.82)
  {\tpcifenglish{half angle: orthogonal $\Rightarrow$ antipodal}%
                {halber Winkel: orthogonal $\Rightarrow$ antipodal}};

% ─────────────────────────────────────────────────────────────
%  (b)  The same arrow, for one spin of an entangled pair
% ─────────────────────────────────────────────────────────────
\begin{axis}[
    tpcplot,
    at={(4.30cm,-1.70cm)}, anchor=south west,
    width=7.0cm, height=4.9cm,
    clip=false,
    xmin=0, xmax=90, ymin=0, ymax=1.42,
    xtick={0,22.5,45,67.5,90},
    xticklabels={$0$,$\pi/8$,$\pi/4$,$3\pi/8$,$\pi/2$},
    ytick={0,0.5,1},
    xlabel={$\alpha$},
    title={(b) \tpcifenglish{a pair: arrow vs.\ entanglement}%
                            {ein Paar: Pfeil gegen Verschränkung}},
    legend style={font=\scriptsize, draw=light-gray,
                  at={(0.5,-0.26)}, anchor=north, legend columns=2},
]

  % The EPR point alpha = pi/4: a neutral reference rule
  \addplot[levellinegrey] coordinates {(45,0) (45,1)};

  % Length of spin 1's Bloch arrow
  \addplot[cgblue, line width=1.0pt, domain=0:90, samples=181]
    {abs(cos(2*x))};
  \addlegendentry{$\tfrac{2}{\hbar}\bigl|\langle\widehat{\mathbf{S}}\otimes\mathbb{I}\rangle\bigr|=|\cos 2\alpha|$}

  % Failure of the product condition c_11 c_22 = c_12 c_21
  \addplot[birrenorange, line width=1.0pt, densely dashed,
           domain=0:90, samples=181] {abs(sin(2*x))};
  \addlegendentry{$2|c_{11}c_{22}-c_{12}c_{21}|=|\sin 2\alpha|$}

  \addplot[only marks, mark=*, mark size=2pt, cgblue, draw=none]
    coordinates {(45,0)};
  \addplot[only marks, mark=*, mark size=2pt, birrenorange, draw=none]
    coordinates {(45,1)};

  % The family the panel is drawn from
  \node[lbl, font=\scriptsize] at (axis cs:45,1.26)
    {$\ket{\psi_\alpha}=\cos\alpha\,\ket{\uparrow\uparrow}
      +\sin\alpha\,\ket{\downarrow\downarrow}$};

  % What the three marked values of alpha are
  \node[lbl, font=\scriptsize, text=birrengray, anchor=north]
    at (axis cs:0,-0.15)  {$\ket{\uparrow\uparrow}$};
  \node[lbl, font=\scriptsize, text=birrengray, anchor=north]
    at (axis cs:45,-0.15) {$\ket{\text{EPR}}$};
  \node[lbl, font=\scriptsize, text=birrengray, anchor=north]
    at (axis cs:90,-0.15) {$\ket{\downarrow\downarrow}$};
\end{axis}

\end{tikzpicture}
\end{center}


Der Pfeil links ist der Erwartungswert $\tfrac{2}{\hbar}\langle\widehat{\mathbf{S}}\rangle$ eines einzelnen Spins. Rechts sieht man ihn schrumpfen, sobald der Spin verschränkt wird: beim EPR-Zustand ist er in den Mittelpunkt der Kugel zurückgewichen -- dort, wo die gestrichelte Kurve ihr Maximum hat.

Solche Zustände heissen \emph{verschränkt}, und sie werden heute im Labor routinemässig erzeugt.
